\section{\bf Introduction}
\label{sec:intro}

% Federated Learning and communication challenges
Federated learning enables collaborative model training on decentralized data by iteratively exchanging model updates, such as gradients. The communication cost of these gradients, however, creates significant overhead.

To address this, gradient compression techniques like quantization, sparsification, and compression are widely used. Yet these costs remain a primary bottleneck that challenges the practical scalability of distributed learning in large training settings.

To this end, more advanced methods improve efficiency by exploiting the inherent correlation between gradients across time and also different nodes, reducing redundant transmissions.

% DSC use in distributed learning
Distributed Source Coding (DSC), based on Slepian-Wolf \cite{SlepianWolf1973} and Wyner-Ziv \cite{1055508} principles, enables workers to independently compress their correlated gradients, and then the parameter server acts as a joint decoder, using gradients from other workers as side information to decompress each update. This exploits inter-node redundancy to reduce communication rate to the joint entropy of the sources, without requiring any direct communication between workers.

% Learned WZ quantizer
Conventional DSC methods rely on rigid algorithms and are not well-suited to large and high-dimensional gradients in deep learning. Recent works \cite{OzyilkanRecoverBinning} propose learned Wyner-Ziv quantizers for Gaussian sources, which could be adapted to gradient distributions and multi-terminal settings. These methods train end-to-end models which include a learned non-uniform quantizer and a dequantizer that jointly processes the incoming code and side information, allowing for lower rates.

This framework can be extended with successive refinement\cite{JoukovskyLayeredWZ}, where a base layer is sent, followed by (optional) residual layers. This offers better rate-distortion tradeoffs and flexibility to adapt to varying bandwidth constraints.

Furthermore, in a multiterminal scenario, we use multiple side information to help the decoding.
The server sequentially decodes updates from each agent, and adds them to a list of side information, allowing for progressively reduced rate for each subsequent worker, improving overall system efficiency.

% Contributions
In this paper, we present a unified framework that combines federated learning with learned successively refined Wyner-Ziv quantizers in a multi-terminal setting. Motivated by the proven bandwidth savings of DSC in FL \cite{ZhangTao2022, Abrahamyan2022LGC} and works on learned WZ quantizers that observed compression gains and scalability, we introduce a fully learned, multi-terminal, end-to-end pipeline specifically designed for gradient compression in distributed learning with special focus on inter-node redundancy.
